% AUTORIDAD DE FUENTE:
% expose_VIII_body.tex, líneas 2171--2203; índices de escaneo 310--311;
% folios de fuente 299--300. La corrección material de la capa francesa
% se conserva reversiblemente mediante \sisi. Los subíndices de fibra
% especial mecanografiados con letra o se representan como cero.

\medskip
\noindent\underline{7. Prolongaciones canónicas de extensiones y
biextensiones por $\underline G_m$.}

\medskip
\noindent 7.0. Recogemos ahora el fruto de nuestros preparativos de las
secciones 5 y 6. Supongamos primero que el esquema base $S$ es un
\underline{rasgo}, es decir, el espectro de un anillo de valoración
discreta; sean $\eta$ su punto genérico, $s$ su punto cerrado y
$k=k(s)$. Sea $P$ un esquema en grupos conmutativo liso sobre $S$, y
escribamos $P_0=P_s$ para su fibra especial. En esta sección solo
consideramos grupos conmutativos; en particular, todas las extensiones
de grupos consideradas son conmutativas.

\medskip
\noindent\underline{Teorema 7.1.} \underline{La notación es la
recordada en 7.0.}

\smallskip
\noindent a) \underline{Supongamos que el grupo de componentes conexas
de la fibra especial de $P$,}
\[
 \Phi=P_0/P_0^{\circ},
\]
\underline{es un grupo de torsión}---como ocurre, en particular, si
$P$ es de tipo finito sobre $S$---. \underline{Entonces el funtor
$E\mapsto E_\eta$ de la categoría de extensiones de $P$ por
$\underline G_{mS}$ a la categoría de extensiones de $P_\eta$ por
$\underline G_{m\eta}$ es plenamente fiel. Supongamos que existe un
entero $n>0$ tal que $n\Phi=0$. Entonces a toda extensión $E_\eta$ de
$P_\eta$ por $\underline G_{m\eta}$ se le asocia un homomorfismo
canónico}
\begin{equation}\tag{7.1.1}
 d(E_\eta):\Phi\longrightarrow(\mathbb Q/\mathbb Z)_k,
\end{equation}
\underline{cuya anulación es necesaria y suficiente para que $E_\eta$
sea isomorfa a la fibra genérica de una extensión de $P$ por
$\underline G_{mS}$. En particular, si $P_0$ es conexo, el funtor
precedente $E\mapsto E_\eta$ es una equivalencia de categorías.}

%%%%%%%%%%  scan idx 311  |  folio 300  %%%%%%%%%%

\smallskip
\noindent b) \underline{Sean $P$ y $Q$ dos esquemas en grupos lisos
sobre $S$, pongamos}
\[
 \Phi=P_0/P_0^{\circ},\qquad \Psi=Q_0/Q_0^{\circ},
\]
\underline{y supongamos que $\Phi$ o $\Psi$ es un grupo de torsión}
(por ejemplo, que $P$ o $Q$ es de tipo finito sobre $S$).
\underline{Entonces el funtor $E\mapsto E_\eta$ de la categoría de
biextensiones de $(P,Q)$ por $\underline G_{mS}$ a la categoría de
biextensiones de $(P_\eta,Q_\eta)$ por $\underline G_{m\eta}$ es
plenamente fiel. Supongamos que existe un entero $n>0$ tal que
$n\Phi=0$ o $n\Psi=0$. Entonces a toda biextensión $E_\eta$ de
$(P_\eta,Q_\eta)$ por $\underline G_{m\eta}$ se le asocia un
emparejamiento canónico}
\begin{equation}\tag{7.1.2}
 d(E_\eta):\Phi\otimes_{\mathbb Z}\Psi
 \longrightarrow(\mathbb Q/\mathbb Z)_k,
\end{equation}
\underline{cuya anulación es necesaria y suficiente para que
$\sisi{E}{E_\eta}$ sea isomorfa a la fibra genérica de una biextensión
$E$ de $(P,Q)$ por
$\sisi{\underline G_{mS}}{\underline G_{mS}}$. En particular, si
$P_0$ o $Q_0$ es conexo, el funtor precedente $E\mapsto E_\eta$ es una
equivalencia de categorías.}
